%% en/sections/03_fazioni.tex — Factions (full EN)

\chapter{The Factions}
\markboth{The Factions}{The Factions}

\lettrine[lines=3]{F}{our} factions move through Palermo in February
1282, unaware or aware of each other, bound by interests that overlap
and contradict. None controls the situation entirely. All, in different
ways, will push 30 March toward what it will become.

\nota[For the Coordinating GM]{
  The factions do not necessarily know each other. Part of the narrative
  tension comes precisely from this opacity: the Conspirators believe
  they control the timing, the Clergy believe they manage the symbols,
  the Barons believe they are negotiating the future, the People believe
  in nothing --- but act. Use the Conjunction Moments to let consequences
  filter through, not causes.
}

\secsep

%% ---- FACTION 1: CONSPIRATORS ----
\section{The Conspirators}
\textcolor{oro}{\cinzel\large The Network of Giovanni da Procida}

\filetto

\textbf{Entry point:} In medias res --- characters are already inside.

\textbf{Session goal:} Ensure the uprising breaks out before Charles's
fleet sails. The plan requires a coordinated signal, but the details
are in the details --- and the details are cracking.

\textbf{Narrative theme:} Trust in a cause does not protect against
betrayal. Those who command from a distance (da Procida, Peter of Aragon,
Michael Palaeologus) have interests that do not coincide with those
who risk their lives in Palermo.

\subsection{Faction Structure}

The Conspirators are a network, not an organisation. Characters know
each other but do not necessarily know all the network's nodes. They
know a plan exists, they know their piece --- but they do not see the
full picture. This opacity is intentional and narratively productive.

\subsection{Faction Score}

\begin{itemize}[nosep]
  \item \textbf{+1} for each baronial contact activated and reliable
  \item \textbf{+1} if the channel to eastern Sicily (Messina) is open
  \item \textbf{+1} if the network's identity has not been compromised
  \item \textbf{+1} if the coordinated signal has been defined and communicated
  \item \textbf{-1} for each network member captured or killed
  \item \textbf{-1} if Angevin authorities are specifically on alert
    for the conspiracy (not generic alertness)
\end{itemize}

\nota[For the Conspirators GM]{
  These characters know more than the others about what is about to happen
  --- but this knowledge is a burden, not an advantage. They are responsible
  for what will occur. If the uprising breaks out the wrong way, it is their
  fault. If it does not break out, it is their fault. The psychological
  pressure is part of the design.
}

\secsep

%% ---- FACTION 2: BARONS ----
\section{The Barons}
\textcolor{oro}{\cinzel\large The Nobility Between Two Fires}

\filetto

\textbf{Entry point:} Recruited --- at session start they are approached
by an emissary (a Conspirator, a friar, or a mysterious intermediary)
with a proposal they cannot ignore.

\textbf{Session goal:} Choose sides --- and survive the choice. The
Sicilian barons have everything to lose: their lands, their titles,
their families. They have everything to gain: the chance to free
themselves from the Angevin yoke. The problem is that both options
are dangerous.

\textbf{Narrative theme:} Privilege is not protection. Having power
means having responsibilities toward those who depend on you --- and
these responsibilities contradict each other.

\subsection{The Barons' Position}

The Sicilian nobility of 1282 was in a structurally unstable position.
Many had collaborated with the Angevins to preserve their lands; some
had obtained new privileges; few had openly opposed them. But none were
satisfied. The Angevin tax system cut into their revenues; French officials
ignored their traditional rights; military requisitions emptied their
vassals.

Alaimo da Lentini, the most important pro-rebel baron of eastern Sicily,
is their natural reference point --- but he is far away, and his position
is as ambiguous as their own.

\subsection{Faction Score}

\begin{itemize}[nosep]
  \item \textbf{+1} if at least three baronial houses have taken position
  \item \textbf{+1} if their own vassals have been prepared and are holding
  \item \textbf{+1} if a secure channel has been established with the
    Conspirators or Clergy
  \item \textbf{+1} if their families are safe
  \item \textbf{-1} for each baronial house that has gone over to the
    Angevin side
  \item \textbf{-1} if a character has been captured or compromised
\end{itemize}

\secsep

%% ---- FACTION 3: CLERGY ----
\section{The Clergy}
\textcolor{oro}{\cinzel\large The Channels of the Church}

\filetto

\textbf{Entry point:} Mixed --- some characters are already aware of
what is being prepared (an informed bishop, a friar used as courier);
others are drawn in by events.

\textbf{Session goal:} Manage the moment without the Church being
compromised --- and at the same time ensure that what is right comes
to pass. The tension between institutional loyalty and individual
conscience is the heart of this table.

\textbf{Narrative theme:} Obedience does not absolve responsibility.
Serving God and serving Rome is not the same thing. And serving your
people is not the same as serving either.

\subsection{The Clergy as Network}

The churches and convents of Palermo are natural communication nodes.
Friars move freely; confessions are inviolable; bell towers are the
highest points in the city. The Clergy has no armies, but something
more precious: it can move symbols, and symbols, at this historical
moment, move people.

The church of Santo Spirito --- outside the walls, site of the Easter
procession --- is the physical centre of the session for this table.

\subsection{Faction Score}

\begin{itemize}[nosep]
  \item \textbf{+1} if the Santo Spirito procession unfolds as planned
    (maximum concentration of people)
  \item \textbf{+1} if communication channels between quarters are
    active and secure
  \item \textbf{+1} if the Church has been kept formally outside the
    conspiracy (plausible deniability)
  \item \textbf{+1} if at least one senior prelate is ready to
    legitimise the uprising after the fact
  \item \textbf{-1} if a clergy member has been arrested or tortured
    by the French
  \item \textbf{-1} if the Angevin authorities have cancelled or
    militarised the procession
\end{itemize}

\nota[For the Clergy GM]{
  This table has access to symbols. Characters can do things no one else
  can: ring the bells at the wrong moment (or right one), open or close
  a church door, hear confession and grant absolution. Use this asymmetry.
  If the Passion layer is active, the Clergy also has the \textit{To Bind
  and to Loose} action (see ``The Passion of the Vespers''): it can forge
  or break other people's Connections by spending Passion.
}

\secsep

%% ---- FACTION 4: PEOPLE ----
\section{The People}
\textcolor{oro}{\cinzel\large The Fuse}

\filetto

\textbf{Entry point:} Dragged in --- characters know nothing of any
conspiracy. They are craftsmen, fishermen, market women, young people
with no prospects. They are recruited or drawn in by someone who trusts
them --- or who needs them.

\textbf{Session goal:} Survive. Protect those you love. And perhaps,
in the process, understand that what is about to happen is not just
a night of blood --- it is something that will change everything.

\textbf{Narrative theme:} History is made by ordinary people, not
heroes. The Vespers did not break out because someone had a plan: they
broke out because ordinary people had reached their limit. This table
plays that limit.

\subsection{Daily Life under the Angevins}

People faction characters live the Angevin pressure concretely and
daily: impossible taxes, requisitions, soldiers searching women, officials
taking what they want. They have no access to great political manoeuvres
--- but they are exactly those who will be in the square on the evening
of 30 March.

\subsection{Faction Score}

\begin{itemize}[nosep]
  \item \textbf{+1} if the popular quarters are informed and mobilised
    (even unknowingly)
  \item \textbf{+1} if a public Angevin abuse has been witnessed in a
    way that can serve as a detonator
  \item \textbf{+1} if neighbourhood networks (the \textit{harat}) are
    ready to act in a coordinated way
  \item \textbf{+1} if the characters have managed to protect someone
    vulnerable (a family, an elderly person, a child) from preliminary
    violence
  \item \textbf{-1} if a key community figure has been eliminated or
    intimidated by the French
  \item \textbf{-1} if the quarters are divided or in conflict with
    each other
\end{itemize}

\nota[For the People GM]{
  This is the most emotionally intense table. Characters have no resources,
  no plans, no protections. They only have each other. Use small details:
  a neighbour's name, the smell of bread, the sound of a night search.
  The big history enters through the servants' entrance.
}

\filettodoppio
